% GPT-6 Astra Pro; effort not exposed; /root/replacement; 22 November 2026.
\section[HD-F1-R28: the surviving product coefficient]{HD-F1-R28: the surviving product character has a nonzero divisor factor and a positive balanced coefficient}
\label{sec:followup-rm-product-character-coefficient}

Retain every actual adjacent object in the precise
R21--R27 class, with its arbitrary ample semistability
polarization $H$ and its determinant twist already
performed. Write the surviving whole characters as
\[
 \operatorname{ch}(T)=a_0v_q+c\,p\overline p,\qquad
 \operatorname{ch}(A)=(a_0+1)v_q+c\,p\overline p,\qquad
 p=r+\beta+sA_1.
 \tag{F1R28.1}
\]
Here $r,s\in F$, $\beta$ is defined over $F$ and
has type $(1,1)$ in the first RM block, and its
conjugate has the stated type in the second block.
These are exactly the field and Hodge conditions
of R26; $c\in\mathbf Q^*$.
For all these actual objects,
\[
       \beta\ne0,\qquad 0<a_0\le2,\qquad
                   a_0=2\ \Longleftrightarrow\ r=0.
 \tag{F1R28.2}
\]
In the three-row common-image branch of R23 the
stronger bound is
\[
                              \frac15<a_0\le2.
 \tag{F1R28.3}
\]
The common-quotient branch and the three-row
common-image branch also force the two divisor
factors to be semidefinite with the exact signs
specified in paragraph~3. Their two ranks agree
and are one or two. These are restrictions on
the whole characters of the original objects;
they do not construct objects at the remaining
coefficient values.

\paragraph{1. The actual effective surface and the nonzero divisor component.}
Use the actual line-hull classes $\ell,m$ and
nonzero effective surface $\mathcal D$ from
(F1R27.19)--(F1R27.22). Thus in the common-image
case they come from $I\subset T$ and $Q=T/I$,
and in the common-quotient case both classes
equal $c_1(L)$ for the actual hull
$T^{**}=L^{\oplus2}$. In both cases,
\[
 \ell H^3=mH^3=0,\qquad
 \ell+m=c_1(T),\qquad
 [\mathcal D]=\tfrac12(\ell^2+m^2)
                       -\operatorname{ch}_2(T),\qquad
                         \mathcal D\ne0.
 \tag{F1R28.4}
\]
All equalities are induced by the same coherent
$K_0$ sequences and the actual leading support
cycles already proved in R27.

Direct multiplication of (F1R28.1) gives
\[
 \begin{split}
 \operatorname{rk}(T)&=a_0+c\,r\overline r=2,\\
 c_1(T)&=c(\overline r\,\beta+r\,\overline\beta),\\
 \operatorname{ch}_2(T)
   &=(-a_0q+c\,\overline r\,s)A_1
        +(-a_0q+c\,r\,\overline s)A_2
        +c\,\beta\overline\beta .
 \end{split}
 \tag{F1R28.5}
\]
In particular $c_1(T)$ is block diagonal.
If $\beta=0$, the degree-two character is in
$\mathbf R A_1+\mathbf R A_2$.
Then (F1R28.4) contradicts the full two-line
surface obstruction (F1R27.4)--(F1R27.5).
Thus $\beta\ne0$.
The scalar-extension involution is invertible,
so $\overline\beta\ne0$ as well, and
$c\beta\overline\beta$ is a nonzero mixed
cohomology tensor.

\paragraph{2. The coefficient bound, with every auxiliary-form map specified.}
When $r=0$, the first equation of (F1R28.5)
already gives $a_0=2$.
Suppose $r\ne0$. Apply the explicit scalar
reparametrization of the same cohomology tensor
\[
 (p,c)\longmapsto
       \left(\frac p r,\,t=c\,r\overline r\right),\qquad
 \frac p r=1+\gamma+\tau A_1,\quad
 \gamma=\frac\beta r,\quad \tau=\frac s r.
 \tag{F1R28.6}
\]
Its identity
$c\,p\overline p=t(p/r)(\overline p/\overline r)$
preserves the entire character. It does not
replace any sheaf or its coherent lift.
Set
\[
       t=2-a_0\ne0,\qquad
       \Gamma=\gamma+\overline\gamma,\qquad
       \eta=\frac{\ell-m}{2}.
 \tag{F1R28.7}
\]
The divisor form $\Gamma$ is real, block
diagonal and rational: it is the
factor-exchange-invariant tensor over $F$.
The first-Chern-class identity is
$\ell+m=t\Gamma$.
Since both actual line-hull classes have
slope zero, one has
\[
                     \Gamma H^3=\eta H^3=0.
 \tag{F1R28.8}
\]

The pure-block squares
$\gamma^2$ and $\overline\gamma^2$ belong to
$\mathbf R A_1$ and $\mathbf R A_2$.
Using
$\gamma\overline\gamma=
 (\Gamma^2-\gamma^2-\overline\gamma^2)/2$
and the exact last line of (F1R28.5), rewrite
the class of the same nonzero effective
surface in (F1R28.4) as
\[
       [\mathcal D]
          =\eta^2-\frac{t a_0}{4}\Gamma^2
                     +uA_1+vA_2
                    \quad\text{for some real }u,v.
 \tag{F1R28.9}
\]
Indeed the coefficient of $\Gamma^2$ is
$t^2/4-t/2=-t a_0/4$ because $t+a_0=2$.
Every equality here is on the particular
invariant representatives of the actual
cohomology classes, as in (F1R27.3).

If $t a_0\le0$, put $\kappa=-t a_0/4\ge0$.
The two explicitly defined real forms
\[
          \lambda_1=\sqrt{\kappa}\Gamma+\eta,
             \qquad
          \lambda_2=\sqrt{\kappa}\Gamma-\eta
 \tag{F1R28.10}
\]
are both $H$-primitive, and their sum is
the block diagonal form $2\sqrt{\kappa}\Gamma$.
Moreover
$\tfrac12(\lambda_1^2+\lambda_2^2)
  =\eta^2+\kappa\Gamma^2$.
Thus (F1R28.9) represents the actual nonzero
effective surface in precisely the prohibited
form (F1R27.5), a contradiction.
The auxiliary real forms need not be
rational or Chern classes: the already
proved positivity lemma applies to these
explicit forms, while its effective cycle
is the unchanged actual $\mathcal D$.
Consequently $t a_0>0$. Together with
$t+a_0=2$ this forces
\[
                             0<a_0<2,\qquad t>0.
 \tag{F1R28.11}
\]
This proves (F1R28.2), including its equality
characterization.

In the three-row common-image case, the
actual source-hull evaluation is
$A^{**}\simeq L^{\oplus3}$ and the target
line-hull classes are $\ell,2\ell$, by
(F1R23.19) and its determinant relation.
Thus $3\ell=t\Gamma$ and
$\eta=-t\Gamma/6$.
For $r\ne0$, (F1R28.9) becomes
\[
       [\mathcal D]=
            \frac{t(2-10a_0)}{36}\Gamma^2
                          +uA_1+vA_2 .
 \tag{F1R28.12}
\]
A nonnegative coefficient would again have
the prohibited two-line form, now with
two identical primitive forms equal to
the nonnegative square root of that
coefficient times $\Gamma$.
The coefficient is therefore negative.
Since $t>0$, this proves $a_0>1/5$.
For $r=0$ the already proved $a_0=2$
gives the same conclusion, proving (F1R28.3).

\paragraph{3. Exact semidefinite divisor factors in the two special branches.}
First suppose $r\ne0$ and use the precise
reparametrization (F1R28.6).
In the common-quotient case, $\ell=m$ and
$\eta=0$, so on a mixed plane the
nonnegative value of (F1R28.9) is
\[
                  -\frac{t a_0}{2}
                    \gamma(a)\overline\gamma(b),
        \qquad a\in T_1,\quad b\in T_2 .
 \tag{F1R28.13}
\]
Here $\gamma(a)$ means the real value of
its Hermitian form on $a$, and likewise
for the second block.
The coefficient is strictly negative
by (F1R28.11).
Thus
$\gamma(a)\overline\gamma(b)\le0$
for every pair of vectors.
Neither form is zero. Fixing a vector
on which one is nonzero proves that
the other has a fixed sign; reversing
the roles proves that both are
semidefinite, with opposite signs.

In the three-row common-image case,
(F1R28.12) gives instead the strictly
negative mixed coefficient
$t(1-5a_0)/9$, by (F1R28.3).
The identical argument proves the same
opposite-sign semidefiniteness of
$\gamma,\overline\gamma$.

When $r=0$, both special branches have
$c_1(T)=0$ and hence their individual
line-hull classes vanish in cohomology:
$2\ell=0$ in the common-quotient case
and $3\ell=0$ in the three-row case.
On mixed planes (F1R28.4)--(F1R28.5)
then give exactly
\[
                    -c\,\beta(a)\overline\beta(b)\ge0
                         \quad\text{for all }a,b.
 \tag{F1R28.14}
\]
Since $c\ne0$ and both forms are nonzero,
they are again semidefinite; their signs
have product $-\operatorname{sgn}(c)$.
This specifies the signs without an
unjustified rescaling by the zero
rank component.

Finally their Hermitian ranks are equal.
In dimension two a nonzero Hermitian
$(1,1)$ form has rank one exactly when
its exterior square is zero.
The coefficient of $\beta^2$ in
the one-dimensional space $F A_1$
is defined over $F$; its conjugate is
the coefficient of $\overline\beta^2$
in $F A_2$. One is zero exactly when
the other is zero. The same holds
after the nonzero scalar change in
(F1R28.6).
Thus the two nonzero semidefinite
factors both have rank one or both
have rank two.

\paragraph{Exact scope and retained framework.}
The surface $\mathcal D$, its coherent
quotient maps, cycle class and invariant
positive representative are exactly
those proved in R22, R25 and R27.
The only additional maps are the scalar
tensor identity (F1R28.6), the explicit
linear combinations of their real
divisor forms, and the degree-two
projection of the whole character.
The sign choice in (F1R28.10) is
essential: the sum of its two forms
is block diagonal even when $\eta$
has off-block components.
No primitive-form replacement of an
actual sheaf is asserted.
The published positivity input and its
actual Haar comparison remain exactly
the Debarre--Ein--Lazarsfeld--Voisin
maps read and proved at their used
scope in R27.

The coefficient bounds and the
nonzero mixed divisor tensor apply
to both factorization branches.
Semidefiniteness has been proved
here only for the common-quotient
and three-row common-image branches.
The two-row common-image branch
retains its actual nonsplit source,
both necessary surface defects,
and the original connecting class.
None of the surviving cases is
asserted to have an actual weakly
semiregular object.
